% Generated draft: requires mathematical and source review.
\documentclass[11pt]{article}
\usepackage[margin=25mm]{geometry}
\usepackage{fontspec}
\setmainfont{Latin Modern Roman}
\usepackage{amsmath,amssymb,mathtools,unicode-math}
\setmathfont{Latin Modern Math}
\usepackage{enumitem,needspace,xurl,hyperref}
\setlength{\parindent}{0pt}
\setlength{\parskip}{5pt}
\setlength{\emergencystretch}{3em}
\newcommand{\N}{\mathbb{N}}
\newcommand{\Z}{\mathbb{Z}}
\newcommand{\Q}{\mathbb{Q}}
\newcommand{\R}{\mathbb{R}}
\newcommand{\C}{\mathbb{C}}
\newcommand{\F}{\mathbb{F}}
\newcommand{\A}{\mathbb{A}}
\newcommand{\PP}{\mathbb{P}}
\newcommand{\E}{\mathbb{E}}
\newcommand{\eps}{\varepsilon}
\newcommand{\dd}{\,\mathrm{d}}
\begin{document}
\section*{UnsolvedMath ID: 1353}
\textbf{Source problem code:} catalogue entry
\subsection{Definitions and mathematical statement}
The supplied question is attributed to W. Hugh Woodin. Its hypothesis is that there is a proper class of Woodin cardinals; equivalently, the Woodin cardinals are unbounded in the ordinals:
\[
\forall\alpha\in\operatorname{Ord}\ \exists\delta\in\operatorname{Ord}:\quad
\alpha<\delta\ \text{ and }\ \delta\text{ is a Woodin cardinal}.
\]
Here a Woodin cardinal is understood in its usual large-cardinal sense. The requested conclusion is an analogue of completeness for $\Omega$-logic: every $\Omega$-semantically valid sentence should admit an $\Omega$-proof, with the appropriate theory parameters and conventions for that logic.

To display the intended direction without inventing missing definitions, write $T\models_{\Omega}\varphi$ for the intended $\Omega$-semantic consequence relation and $T\vdash_{\Omega}\varphi$ for the intended $\Omega$-provability relation. The source's target has the schematic form
\[
\bigl[\text{there is a proper class of Woodin cardinals}\bigr]
\ \Longrightarrow\ 
\bigl[T\models_{\Omega}\varphi\ \Longrightarrow\ T\vdash_{\Omega}\varphi\bigr]
\]
for the stipulated class of theories $T$ and sentences $\varphi$. This is a schematic description of the proposed completeness implication, not a fully specified theorem: the supplied record does not define either relation, the class of permitted theories, the language, or whether the hypothesis is imposed in the ambient universe, as part of $T$, or both. The word ``appropriate'' in the supplied assessment cannot fix these omitted choices. They must be recovered from a precise source before this implication is treated as a complete formalization.

The background says that $\Omega$-logic uses Woodin cardinals to define semantic validity and describes possible consequences for set-theoretic truth and independence questions. Those motivational claims are not additional quantified conclusions here. The literature triage, checked 2026-08-17, labels Woodin's $\Omega$-conjecture open and reports conditional $\Omega$-logic results and the Strong $\Omega$ Conjecture as a framework for generic-absoluteness consequences. It does not identify those conditional results precisely or define the Strong $\Omega$ Conjecture. The strong variant is therefore not substituted for the completeness question in this record. The triage states that no completeness theorem of the stipulated form was verified; the target remains conditional on the stated large-cardinal hypothesis.
\subsection{Short English statement}
Assuming unboundedly many Woodin cardinals, must every semantically valid statement of $\Omega$-logic have an $\Omega$-proof? The supplied source identifies the conjecture but omits the exact semantic and proof relations and their theory parameters, so the displayed implication remains a marked schema requiring source review.
\subsection{Sources}
\begin{itemize}
\item UnsolvedMath local registry, permanent ID 1353, source code catalogue entry, ``Woodin's $\Omega$-Conjecture'', attributed in the background to W. Hugh Woodin. Dataset attribution: Ulam AI and the UnsolvedMath Contributors; curation and metadata under CC BY 4.0.
\item Stanford Encyclopedia of Philosophy, \emph{The Continuum Hypothesis}, \url{https://plato.stanford.edu/entries/continuum-hypothesis/}. The supplied triage labels this an authoritative secondary source stating that the corresponding $\Omega$-logic completeness theorem is open.
\item The local literature triage is dated 2026-08-17. No primary paper, exact formulation, or definitions of the $\Omega$-logic relations are supplied in the record.
\end{itemize}
\end{document}
